Au 1\up{er} janvier 2025, une réserve d'eau naturelle contient $\np{3000}\text{ m}^{3}$ d'eau.

Une étude indique que la réserve perd chaque année 10\,\% du volume d'eau qu'elle avait en début d'année.

\medskip

On modélise l'évolution du volume d'eau de cette réserve par une suite $(u_{n})$, où $u_{n}$ désigne la quantité d'eau, exprimée en mètres cubes, contenue dans la réserve au 1\up{er} janvier de l'année $2025+n$.

On a donc $u_{0}=\np{3000}$.

\begin{enumerate}
	\item Montrer que $u_{1}=\np{2700}$.
	\item Pour tout entier naturel $n$, exprimer $u_{n+1}$ en fonction de $u_{n}$.
	\item Quelle est la nature de la suite $(u_{n})$ ? Préciser sa raison.
	\item Dans la feuille de calcul suivante on a calculé les premiers termes de la suite $(u_{n})$. Le format des cellules de la ligne 2 permet d'afficher les nombres arrondis à l'unité.

	\begin{center}
	\begin{tblr}[T]
		{hlines, vlines, columns={c},
		cells={font=\sffamily},
		row{1}={bg=gray9},
		column{1}={bg=gray9},
		cell{1}{1}={gray9,cmd={\tikz[baseline={(0,0.05)}]{\clip(0,0) rectangle (0.6,0.5) ;\fill[gray!50] (0,0)--(0.5,0.5)--(0.5,0);}}} }
		   & A & B & C & D & E & F & G & H & I \\
		1  & $n$ & 0 & 1 & 2 & 3 & 4 & 5 & 6 & 7 \\
		2  & $u_n$ & \np{3000} & \np{2700} & \np{2430} & \np{2187} & \np{1968} & \np{1771} & \np{1594} & \np{1435} \\
	\end{tblr}
	\end{center}

	Quelle formule a été saisie en cellule C2 afin d'obtenir, par glisser-coller vers la droite, les premiers termes de la suite $(u_{n})$ ?

	\item Pour assurer la sécurité de la lutte contre les incendies on estime que la réserve doit contenir au moins $\np{2000}\text{ m}^{3}$.

	D'après le tableau précédent, à partir de quelle année la quantité d'eau dans la réserve deviendra insuffisante ?
\end{enumerate}
